\documentclass[./main.tex]{subfiles}
\begin{document}

% Slide # 7
\begin{frame}[t, label=slide07]
        % Title
        \frametitle{Theorem n. 1 (existence and uniqueness)}

        % Body
        \footnotesize

        \begin{theorem}[Picard-Lindel{\"o}f]\label{thm:picard_lindelof}
             Let $\dot{x}=f(x,t)$ with $f:\,D\subseteq \mathbb{R}^{n}\times \mathbb{R}\,\to \mathbb{R}^{n}$ continuous in $t$ and Lipschitz in $x$ and $x(t_{0})=x_{0}$ being the initial datum. Then there exist a $\delta>0$ s.t. the solution $\phi(x_{0},t_{0})$ exists and is unique on $\mathbb{B}_{\delta}(t_{0})$. 
        \end{theorem}
        
        \begin{columns}[T]
                
                \column{0.5\textwidth}

                In other words, if there are two solutions $\phi_{1}$ and $\phi_{2}$ defined on the same time interval $\mathbb{B}_{\delta}(t_{0})$ stemming out from the same initial condition $x_{0}$, then $\phi_{1}=\phi_{2}$.

                \vspace{-0.5mm}
                \begin{figure}[H]
                        \centering 
                        \includegraphics<2->[keepaspectratio, width=\textwidth]{./figures/slide_07_1.png}%
                \end{figure}

                \hfill

                \column{0.5\textwidth}

                \vspace{-5mm}
                \begin{figure}[H]
                        \centering 
                        \includegraphics<2->[keepaspectratio, width=\textwidth]{./figures/slide_07_2.png}%
                \end{figure}

       \end{columns}



\end{frame}

\end{document}
